\chapter{Hermitian line twists}

\begin{lemma}[Semistability under a Hermitian line twist]
\label{lem:line-twist}
\uses{lem:degree-calculus}
\source{bost-specht-manuscript:page-0012,bost-specht-manuscript:page-0013}
\notready
Let $K$ be any number field, $L$ a rank-one Hermitian projective
$\Ocal_K$-module, and $E$ a positive-rank Hermitian projective
$\Ocal_K$-module. Tensoring by $L$ preserves semistability and shifts
every positive-rank submodule slope, as well as the ambient slope,
by $\degn(L)$. In particular, if $E$ is semistable, both actual
Hilbert tensor products $L\otimes E$ and $E\otimes L$ are semistable.
No principal-ideal or freeness hypothesis on $L$ is imposed.
\end{lemma}
\begin{proof}
An invertible module $L$ gives inverse correspondences
$B\mapsto L\otimes B$ and $C\mapsto L^\vee\otimes C$ between
submodules of $E$ and of $L\otimes E$. Local trivializations show
that ranks and saturation are preserved. At each archimedean place,
the line tensor factor multiplies every vector norm by the norm of a
line generator; hence these correspondences preserve the induced
metric up to that same line twist. The degree formula shifts every
submodule slope, as well as the ambient slope, by $\degn(L)$, so
their comparisons are unchanged.
The tensor symmetry $E\otimes L\simeq L\otimes E$ preserves the
Hilbert metric: this follows on pure tensors from the product formula
for the inner product, and then on all tensors by sesquilinearity.
It therefore transfers the same assertion to the other tensor order.
\end{proof}
